% GENERATED FILE. Edit canonical lessons, then run npm run generate:books.
% canonical-source-hash: fnv1a64:a24f19b861bdb612
% canonical-lessons: ZH-C511-qinren, ZH-W511-qin-kin, ZH-C511-fuqin, ZH-C511-muqin, ZH-C511-qinsheng, ZH-C511-qinzi

\chapter{Close Family}
\label{ch:zh-close-family}
\hlchaptermodality{511}

\begin{chapteropening}
\textbf{By the end of this chapter:} \emph{I can talk about my nearest and dearest, my father and mother, children by birth and doing things myself.}

Complete the last lesson of chapter 511: I can talk about my nearest and dearest, my father and mother, children by birth and doing things myself.
\end{chapteropening}

\section[\texorpdfstring{\zh{亲人} (qīnrén)}{亲人 (qīnrén)}]{\zh{亲人} (qīnrén) — one's close family; one's nearest and dearest}

\label{lesson:ZH-C511-qinren}



\begin{quote}
Before the new one: say the Mandarin for to dote on, then the Mandarin for to soothe a child.
\end{quote}

\subsection*{You'll want to know: \zh{亲人}}
\textbf{\zh{亲人}} — \emph{qīnrén} — \textquotedblleft{}one's close family; one's nearest and dearest\textquotedblright{}.

\begin{quote}
\textbf{\zh{在国外的时候我最想我的亲人}} \emph{zài guówài de shíhou wǒ zuì xiǎng wǒ de qīnrén} — When I was abroad, I missed my family most of all.
\end{quote}

The new character, in pieces: \zh{亲} — nine strokes in two pieces: on top the top of \zh{意} \emph{yì}, five strokes; under it four strokes, a bar, a vertical that hooks at the foot, a stroke falling left and a dot. The next lesson puts it in your hand.

\begin{cousinweb}
\textbf{\zh{亲人}} (\emph{qīnrén}) is built from \zh{亲} 'close' and \zh{人} 'a person', both inherited from Middle Chinese, \emph{tshin} and \emph{nyin} (Guangyun, in Baxter's transcription), which glosses the first 'to love' and 'near'.
\end{cousinweb}

\begin{grammarlens}[title={What you've built}]
One more word for saying how it was.
\end{grammarlens}

\subsection*{Your turn}
\begin{itemize}
\raggedright
  \item \emph{Say it:} \emph{qīnrén}
  \item \emph{Say it:} \emph{qīnrén}, once more
  \item \emph{Say it:} \emph{zài guówài de shíhou wǒ zuì xiǎng wǒ de qīnrén} — When I was abroad, I missed my family most of all.
  \item \emph{Recall:} say the Mandarin for to dote on, then the Mandarin for to soothe a child, then say \emph{qīnrén} again
\end{itemize}

\begin{tcolorbox}[breakable,colback=teal!4,colframe=teal!35!black,title={Before you move on}]
What does \zh{亲人} mean? (\textquotedblleft{}One's close family; one's nearest and dearest\textquotedblright{}.) Say it once more.
\end{tcolorbox}
\section[\texorpdfstring{\zh{亲} (qīn)}{亲 (qīn)}]{\zh{亲} — the top of yì, then four strokes}

\label{lesson:ZH-W511-qin-kin}



\begin{quote}
The word first, then the character inside it.

\begin{itemize}
\raggedright
  \item \emph{Recall:} say \emph{qīnrén} — \textbf{R1}, one lesson back, before it decays
\end{itemize}

\begin{quote}
\zh{意} \zh{人}
\end{quote}

\emph{Read it:} each one aloud, with nothing said first — \emph{yì}, \emph{rén}
\end{quote}

\subsection*{Script — the top of yì, then four strokes}
\hlblockfigure{\detokenize{figures/ZH-W511-qin-kin-filmstrip.pdf}}{How \zh{亲} is written, stroke by stroke}

\begin{quote}
\zh{亲}
\end{quote}

\textbf{Two pieces, top before bottom.}

\textbf{The top, five strokes}, as on top of \zh{意}: a dot, a bar, a dot on the left and a short stroke falling to the left on the right, then a long bar.

\textbf{The bottom, four strokes}: a bar, a vertical down the middle that hooks to the left at the foot, a short stroke falling to the left, and a dot falling to the right.

5 + 4 = nine.

\subsection*{Your turn}
\begin{itemize}
\raggedright
  \item \emph{Write it:} \zh{亲}, all nine, with the model in front of you
  \item \emph{Cover:} write it again from memory, and count the pieces, not the strokes
\end{itemize}

\begin{tcolorbox}[breakable,colback=teal!4,colframe=teal!35!black,title={Before you move on}]
How many pieces does \zh{亲} have? (\textbf{Two}.) Which character you know has the same top? (\textbf{\zh{意}}.) How does the middle vertical end? (\textbf{With a hook to the left}.)

Source: \href{https://raw.githubusercontent.com/chanind/hanzi-writer-data/68d10a4b21150cae5e1ebbd223eed289cf32d90c/data/\%E4\%BA\%B2.json}{Hanzi Writer Data \zh{亲}, pinned stroke paths and medians}.
\end{tcolorbox}
\section[\texorpdfstring{\zh{父亲} (fùqīn)}{父亲 (fùqīn)}]{\zh{父亲} (fùqīn) — father (the formal word)}

\label{lesson:ZH-C511-fuqin}



\begin{quote}
Before the new one: say the Mandarin for to soothe a child, then the Mandarin for one's nearest and dearest.
\end{quote}

\subsection*{You'll want to know: \zh{父亲}}
\textbf{\zh{父亲}} — \emph{fùqīn} — \textquotedblleft{}father (the formal word)\textquotedblright{}.

\begin{quote}
\textbf{\zh{我父亲在大学教了三十年汉语}} \emph{wǒ fùqīn zài dàxué jiāo le sānshí nián Hànyǔ} — My father taught Chinese at a university for thirty years.
\end{quote}

Many speakers make the second syllable light: fùqin.

\begin{cousinweb}
\textbf{\zh{父亲}} (\emph{fùqīn}) is built from \zh{父} 'a father' and \zh{亲} 'close', both inherited from Middle Chinese, \emph{bjuX} and \emph{tshin} (Guangyun, in Baxter's transcription).
\end{cousinweb}

\begin{grammarlens}[title={What you've built}]
One more word for saying how it was.
\end{grammarlens}

\subsection*{Your turn}
\begin{itemize}
\raggedright
  \item \emph{Say it:} \emph{fùqīn}
  \item \emph{Say it:} \emph{fùqīn}, once more
  \item \emph{Say it:} \emph{wǒ fùqīn zài dàxué jiāo le sānshí nián Hànyǔ} — My father taught Chinese at a university for thirty years.
  \item \emph{Recall:} say the Mandarin for to soothe a child, then the Mandarin for one's nearest and dearest, then say \emph{fùqīn} again
\end{itemize}

\begin{tcolorbox}[breakable,colback=teal!4,colframe=teal!35!black,title={Before you move on}]
What does \zh{父亲} mean? (\textquotedblleft{}Father (the formal word)\textquotedblright{}.) Say it once more.
\end{tcolorbox}
\section[\texorpdfstring{\zh{母亲} (mǔqīn)}{母亲 (mǔqīn)}]{\zh{母亲} (mǔqīn) — mother (the formal word)}

\label{lesson:ZH-C511-muqin}



\begin{quote}
Before the new one: say the Mandarin for one's nearest and dearest, then the Mandarin for father (the formal word).
\end{quote}

\subsection*{You'll want to know: \zh{母亲}}
\textbf{\zh{母亲}} — \emph{mǔqīn} — \textquotedblleft{}mother (the formal word)\textquotedblright{}.

\begin{quote}
\textbf{\zh{我母亲比我父亲大两岁}} \emph{wǒ mǔqīn bǐ wǒ fùqīn dà liǎng suì} — My mother is two years older than my father.
\end{quote}

Many speakers make the second syllable light: mǔqin.

\begin{cousinweb}
\textbf{\zh{母亲}} (\emph{mǔqīn}) is built from \zh{母} 'a mother' and \zh{亲} 'close', both inherited from Middle Chinese, \emph{muwX} and \emph{tshin} (Guangyun, in Baxter's transcription).
\end{cousinweb}

\begin{grammarlens}[title={What you've built}]
One more word for saying how it was.
\end{grammarlens}

\subsection*{Your turn}
\begin{itemize}
\raggedright
  \item \emph{Say it:} \emph{mǔqīn}
  \item \emph{Say it:} \emph{mǔqīn}, once more
  \item \emph{Say it:} \emph{wǒ mǔqīn bǐ wǒ fùqīn dà liǎng suì} — My mother is two years older than my father.
  \item \emph{Recall:} say the Mandarin for one's nearest and dearest, then the Mandarin for father (the formal word), then say \emph{mǔqīn} again
\end{itemize}

\begin{tcolorbox}[breakable,colback=teal!4,colframe=teal!35!black,title={Before you move on}]
What does \zh{母亲} mean? (\textquotedblleft{}Mother (the formal word)\textquotedblright{}.) Say it once more.
\end{tcolorbox}
\section[\texorpdfstring{\zh{亲生} (qīnshēng)}{亲生 (qīnshēng)}]{\zh{亲生} (qīnshēng) — one's own (by birth)}

\label{lesson:ZH-C511-qinsheng}



\begin{quote}
Before the new one: say the Mandarin for father (the formal word), then the Mandarin for mother (the formal word).
\end{quote}

\subsection*{You'll want to know: \zh{亲生}}
\textbf{\zh{亲生}} — \emph{qīnshēng} — \textquotedblleft{}one's own (by birth)\textquotedblright{}.

\begin{quote}
\emph{tāmen bú shì qīnshēng xiōngdì} — They aren't brothers by birth.
\end{quote}

\begin{cousinweb}
\textbf{\zh{亲生}} (\emph{qīnshēng}) is built from \zh{亲} 'close' and \zh{生} 'to be born', both inherited from Middle Chinese, \emph{tshin} and \emph{srjaeng} (Guangyun, in Baxter's transcription), which glosses the second 'to grow'.
\end{cousinweb}

\begin{grammarlens}[title={What you've built}]
One more word for saying how it was.
\end{grammarlens}

\subsection*{Your turn}
\begin{itemize}
\raggedright
  \item \emph{Say it:} \emph{qīnshēng}
  \item \emph{Say it:} \emph{qīnshēng}, once more
  \item \emph{Say it:} \emph{tāmen bú shì qīnshēng xiōngdì} — They aren't brothers by birth.
  \item \emph{Recall:} say the Mandarin for father (the formal word), then the Mandarin for mother (the formal word), then say \emph{qīnshēng} again
\end{itemize}

\begin{tcolorbox}[breakable,colback=teal!4,colframe=teal!35!black,title={Before you move on}]
What does \zh{亲生} mean? (\textquotedblleft{}One's own (by birth)\textquotedblright{}.) Say it once more.
\end{tcolorbox}
\section[\texorpdfstring{\zh{亲自} (qīnzì)}{亲自 (qīnzì)}]{\zh{亲自} (qīnzì) — personally (doing it oneself)}

\label{lesson:ZH-C511-qinzi}



\begin{quote}
Before the new one: say the Mandarin for mother (the formal word), then the Mandarin for one's own (by birth).
\end{quote}

\subsection*{You'll want to know: \zh{亲自}}
\textbf{\zh{亲自}} — \emph{qīnzì} — \textquotedblleft{}personally (doing it oneself)\textquotedblright{}.

\begin{quote}
\textbf{\zh{妈妈亲自来学校看我}} \emph{māma qīnzì lái xuéxiào kàn wǒ} — Mum came to school to see me herself.
\end{quote}

\begin{cousinweb}
\textbf{\zh{亲自}} (\emph{qīnzì}) is built from \zh{亲} 'close' and \zh{自} 'self', both inherited from Middle Chinese, \emph{tshin} and \emph{dzijH} (Guangyun, in Baxter's transcription), which glosses the second 'from'; the sense 'self' is not given there.
\end{cousinweb}

\begin{grammarlens}[title={What you've built}]
One more word for saying how it was.
\end{grammarlens}

\subsection*{Your turn}
\begin{itemize}
\raggedright
  \item \emph{Say it:} \emph{qīnzì}
  \item \emph{Say it:} \emph{qīnzì}, once more
  \item \emph{Say it:} \emph{māma qīnzì lái xuéxiào kàn wǒ} — Mum came to school to see me herself.
  \item \emph{Recall:} say the Mandarin for mother (the formal word), then the Mandarin for one's own (by birth), then say \emph{qīnzì} again
\end{itemize}

\begin{tcolorbox}[breakable,colback=teal!4,colframe=teal!35!black,title={Before you move on}]
What does \zh{亲自} mean? (\textquotedblleft{}Personally (doing it oneself)\textquotedblright{}.) Say it once more.
\end{tcolorbox}
